% CompletenessThermalFields20261024: independent negative-order normalization.
这里 $l,m$ 均为整数，先取 $0\le m\le l$，并使用含 Condon--Shortley 相位 $(-1)^m$ 的约定。
为避免用待证关系定义负阶，\emph{独立规定} $P_l^{-m}$ 是同一连带 Legendre 方程在两端正则的解，且
\begin{equation}
 \lim_{x\to1^-}\frac{P_l^{-m}(x)}{(1-x^2)^{m/2}}=\frac1{2^m m!}.\tag{A3.N}
 \label{advA3:eq:complete-negative-norm}
\end{equation}
此规范给出函数的幅度和相位，不引用正、负阶之间的阶乘比。
当 $m=0$ 时它就是 $P_l(1)=1$。
端点指标为 $\pm m/2$；$m>0$ 时取正指标的正则支，$m=0$ 时舍去对数发散支。
下题将由这个独立规范证明负阶关系，同时说明这样的解存在且唯一。
